\documentclass[12pt]{article}
\usepackage{amsmath}
\usepackage{geometry}
\geometry{letterpaper, margin=1in}
\usepackage{fancyhdr}
\usepackage{setspace}
\usepackage{hyperref}

\pagestyle{fancy}
\fancyhf{}
\rhead{33-120 Science and Science Fiction}
\lhead{Sci-Fi Problem 2 (25 points)}
\setlength{\headheight}{15pt}

\begin{document}

\vspace{1cm}

\noindent \textbf{Concept}: \textit{Relativistic time dilation}—Moving clocks
run slow, compared to clocks at rest.

\[
\Delta t = \frac{\Delta t'}{\sqrt{1 - \frac{v^2}{c^2}}}
\]

\noindent In the equation above, $\Delta t'$ is the time interval as perceived
in the moving frame of reference, $\Delta t$ is the time interval as perceived
in the rest frame (not moving), $v$ is the speed of the moving reference frame
relative to the rest frame, and $c$ is the speed of light: $c = 2.998 \times
10^8$ m/s.

\vspace{0.5cm}

\noindent You will receive 5 points for an on time submission. The final 20
points are assigned below. In order to get full credit for the problems below,
be sure to:
\begin{itemize}
    \item explicitly state all assumptions made that are not given in the
    problem.
    \vspace{-10pt}
    \item show all of the steps in each calculation.
    \vspace{-10pt}
    \item include proper physical units for all results.
\end{itemize}

\vspace{10pt}
\noindent \textbf{Problem 1. (20 points)} Imagine human space flight from Earth
to Mars in the following two stages.

\begin{enumerate}
    \item (4 pts) Go to NASA's website and look up information about the speed
    of the International Space Station (ISS) in its orbit around the Earth.
    Express this speed in units of meters per second.

    \item (8 pts) Use Einstein's equation for time dilation to calculate how
    much slower clocks on the ISS run, compared to clocks at rest on the surface
    of the Earth. You can do this by calculating the ratio $\frac{\Delta
    t'}{\Delta t}$ (time elapsed on the ISS divided by time elapsed on Earth).
    \textbf{Note:} The answer should come out very close to 1, but the goal of
    this problem is to discovery how different from 1 it really is. Be sure to
    do this with a computer or a calculator that can show enough significant
    digits. If you round off to 1, you will not receive points for this part. 

    \item (8 pts) Suppose one of the astronauts on the ISS has an identical twin
    of exactly the same age. If the astronaut spends a full year on the ISS
    (i.e., 1 year as experienced in the moving frame of reference of the ISS),
    how much older will the twin be who remained on Earth the whole time? 
\end{enumerate}


\newpage
\vspace{10pt}
\noindent \textbf{Problem 1. (20 points)} Imagine human space flight from Earth
to Mars in the following two stages.

\begin{enumerate}
    \item (4 pts) Go to NASA's website and look up information about the speed
    of the International Space Station (ISS) in its orbit around the Earth.
    Express this speed in units of meters per second.

    \item Answer (2pts for trying; 2 pts for right answer at order of magnitude; -2 for no source):
    \href{https://www.nasa.gov/international-space-station/space-station-facts-and-figures/}{https://www.nasa.gov/international-space-station/space-station-facts-and-figures/}:
    5 miles per second converts to 8047 m/s 


    \item (8 pts) Use Einstein's equation for time dilation to calculate how
    much slower clocks on the ISS run, compared to clocks at rest on the surface
    of the Earth. You can do this by calculating the ratio $\frac{\Delta
    t'}{\Delta t}$ (time elapsed on the ISS divided by time elapsed on Earth).
    \textbf{Note:} The answer should come out very close to 1, but the goal of
    this problem is to discovery how different from 1 it really is. Be sure to
    do this with a computer or a calculator that can show enough significant
    digits. If you round off to 1, you will not receive points for this part. 

    \item Answer (2 pts for trying; 4 pts for correct algebraic equation; 2 pts for correct number plug in.): \[
        \frac{\Delta t'}{\Delta t} = \sqrt{1 - \frac{v^2}{c^2}} = \sqrt{1-(8047 \rm{m/s})^2/(3\times10^8 \rm{m/s})^2 } = 0.99999999961
           \]

    \item (8 pts) Suppose one of the astronauts on the ISS has an identical twin
    of exactly the same age. If the astronaut spends a full year on the ISS
    (i.e., 1 year as experienced in the moving frame of reference of the ISS),
    how much older will the twin be who remained on Earth the whole time? 

    \item Answer (2 pts for trying; 4 pts for correct algebraic equation; 2 pts for correct number plug in.): \[\Delta t = \frac{\Delta t'}{\sqrt{1 - \frac{v^2}{c^2}}} = \frac{1\,\rm{yr}}{\sqrt{1-(8047 \rm{m/s})^2/(3\times10^8 \rm{m/s})^2 }} = 1 \rm{yr} /0.99999999961 = 1.00000000039 \rm{yr} \] So the twin is 0.00000000039 yr = 0.0123 seconds older. 
\end{enumerate}

\end{document}
